
\section{Empirical Verification}

\subsection{Design}

The empirical exercise has two layers. The equity-option layer uses monthly rolling option
buckets on AAPL, AMZN, GOOGL, JPM, META, MSFT, NVDA, and TSLA. Each date-bucket
representative is the liquid contract with the highest volume and acceptable relative
spread in one of four buckets: ATM, near put, near call, and put wing. The VIX-option
layer is built separately from raw Databento OPRA VIX chain shards. It preserves contract
identity, parses OSI expiry/right/strike, aligns each contract to the nearest eligible VX
futures forward, computes Black--76 implied volatility and Greeks only after monthly
representative selection, and then constructs an expiry holding ledger.

The train window ends on 2020-12-31; all main performance statistics are reported for the
subsequent 60 monthly test observations. Expected returns, covariance estimates, residuals,
betas, state variables, borrow proxies, and cost inputs are read only from training-window
or decision-date information. The gross return is the clean option-only Markowitz object:
selected option payoff or mark return times signed premium/NAV weight. The net research
return subtracts a conservative implementation layer for spreads, explicit fees, slippage,
borrow, capacity, simulated margin funding, and short-option assignment-risk flags. This
net layer is a research simulation, not a broker-executed fill model.

VIX settlement is handled as a gating variable rather than a footnote. A listed VIX option
expiry backtest requires exact VRO/SOQ settlement. The runner accepts exact local VRO/SOQ
files through the configured settlement path. If exact settlement is absent for any VIX
expiry row, the row remains a VIX-close settlement proxy and VIX-enabled results are marked
diagnostic rather than headline-grade. Table~\ref{tab:vixsettlement} is therefore part of
the claim, not just data documentation.

\begin{table}[H]
\centering
\resizebox{0.94\textwidth}{!}{\begin{tabular}{ll}
\toprule
Item & Value \\
\midrule
Raw OPRA-derived equity option rows & 160,315 \\
Raw VIX option rows after filters & 103,650 \\
Monthly snapshot dates & 129 \\
Training dates & 69 \\
Test dates & 60 \\
Primary equity underlyings & AAPL, AMZN, GOOGL, JPM, META, MSFT, NVDA, TSLA \\
VIX option treatment & VX-forward Black-76 Greeks; VIX/VVIX state only \\
VIX settlement source & exact VRO/SOQ (536 rows) \\
VIX headline status & exact VRO/SOQ settlement (excluded expiries disclosed in Section 2) \\
Post-cost layers & executable-cost scenarios (mid, half, full spread) and conservative stress-cost layer \\
Broker execution status & not broker-executed live evidence \\
Rolling option bucket assets & 54 \\
Train/test split & through 2020-12-31 / after 2020-12-31 \\
\bottomrule
\end{tabular}
}
\caption{Data, timing, settlement, and post-cost construction summary generated by the
empirical runner.}
\label{tab:empirical_data}
\end{table}

\begin{table}[H]
\centering
\resizebox{0.62\textwidth}{!}{\begin{tabular}{lll}
\toprule
Settlement source & Rows & Headline eligible \\
\midrule
vro\_soq\_exact & 536 & yes \\
\bottomrule
\end{tabular}
}
\caption{VIX settlement coverage. VIX rows are headline-grade only when every VIX expiry
row used in headline results has \texttt{vro\_soq\_exact} settlement.}
\label{tab:vixsettlement}
\end{table}

\subsection{Performance and implementation robustness}

\begin{table}[H]
\centering
\resizebox{\textwidth}{!}{\begin{tabular}{llrrrrrrr}
\toprule
Return basis & Strategy & Ann. return & Ann. vol & Sharpe & Sortino & Calmar & Omega & Info. ratio \\
\midrule
Gross before costs & Equity-option Greek Markowitz & 0.894 & 1.062 & 0.842 & 1.770 & 1.020 & 1.888 & 0.598 \\
Gross before costs & Greek Markowitz + VIX & 1.246 & 0.907 & 1.374 & 3.385 & 2.034 & 2.784 & 1.200 \\
Gross before costs & Beta/delta-neutral + VIX & 1.172 & 0.829 & 1.414 & 3.410 & 2.057 & 2.948 & 1.147 \\
Gross before costs & Cost-aware Sortino + VIX & 5.020 & 4.949 & 1.014 & 3.349 & 4.993 & 2.471 & 0.971 \\
Gross before costs & Equal premium & -0.056 & 1.878 & -0.030 & -0.055 & -0.056 & 0.977 & -0.120 \\
Gross before costs & Equal risk & 0.104 & 1.275 & 0.082 & 0.146 & 0.107 & 1.064 & -0.058 \\
Gross before costs & VIX hedge sleeve & -9.285 & 1.444 & -6.429 & -3.089 &  & 0.040 & -6.132 \\
Gross before costs & Delta-matched equities & 1.095 & 0.931 & 1.176 & 2.317 & 1.225 & 2.367 & 1.041 \\
Gross before costs & Underlying Markowitz & 0.179 & 0.196 & 0.909 & 1.856 & 0.590 & 1.929 & 0.072 \\
Post-cost research & Equity-option Greek Markowitz & -7.900 & 5.476 & -1.443 & -1.347 & -6.357 & 0.098 & -1.545 \\
Post-cost research & Greek Markowitz + VIX & -6.079 & 4.827 & -1.259 & -1.198 & -5.136 & 0.116 & -1.345 \\
Post-cost research & Beta/delta-neutral + VIX & -3.896 & 2.841 & -1.371 & -1.300 & -3.895 & 0.151 & -1.508 \\
Post-cost research & Cost-aware Sortino + VIX & -157.425 & 148.976 & -1.057 & -1.019 & -1.985 & 0.016 & -1.003 \\
Post-cost research & Equal premium & -3.165 & 2.378 & -1.331 & -1.502 & -3.051 & 0.359 & -1.372 \\
Post-cost research & Equal risk & -1.276 & 1.367 & -0.934 & -1.199 & -1.276 & 0.505 & -1.082 \\
Post-cost research & VIX hedge sleeve & -14.364 & 4.209 & -3.413 & -2.445 & -12.391 & 0.016 & -3.357 \\
Post-cost research & Delta-matched equities & 1.095 & 0.931 & 1.176 & 2.317 & 1.225 & 2.367 & 1.041 \\
Post-cost research & Underlying Markowitz & 0.179 & 0.196 & 0.909 & 1.856 & 0.590 & 1.929 & 0.072 \\
\bottomrule
\end{tabular}
}
\caption{Gross and post-cost out-of-sample performance. Gross returns test the option-only
Markowitz object; post-cost returns subtract the conservative research implementation
layer. Headline eligibility is downgraded when exact VIX settlement is incomplete.}
\label{tab:performance}
\end{table}

\begin{table}[H]
\centering
\resizebox{\textwidth}{!}{\begin{tabular}{lrrrrrrrr}
\toprule
Strategy & Mean monthly cost & Ann. cost drag & Spread share & Fee share & Mean margin/NAV & Mean stress margin/NAV & Mean assignment notional/NAV & Capacity penalized share \\
\midrule
Beta/delta-neutral + VIX & 0.422 & 5.068 & 0.099 & 0.012 & 4.029 & 0.235 & 0.000 & 0.273 \\
Equal premium & 0.259 & 3.109 & 0.115 & 0.025 & 0.953 & 0.238 & 0.000 & 0.381 \\
Equal risk & 0.115 & 1.381 & 0.213 & 0.035 & 0.920 & 0.230 & 0.000 & 0.347 \\
Equity-option Greek Markowitz & 0.733 & 8.794 & 0.021 & 0.003 & 4.610 & 0.224 & 0.039 & 0.335 \\
Greek Markowitz + VIX & 0.610 & 7.325 & 0.069 & 0.008 & 3.526 & 0.234 & 0.000 & 0.290 \\
VIX hedge sleeve & 0.430 & 5.165 & 0.347 & 0.046 & 0.978 & 0.244 & 0.000 & 0.847 \\
\bottomrule
\end{tabular}
}
\caption{Research implementation diagnostics. Costs include spread crossing, explicit
fees, slippage, borrow, capacity penalties, simulated margin funding, and assignment-risk
penalties where flagged. These are not live broker fills.}
\label{tab:costs}
\end{table}

Table~\ref{tab:performance} separates the mathematical portfolio object from the economic
implementation test. The gross option strategies can look attractive before costs. The
post-cost research simulation is deliberately stricter: wide option spreads, capacity
penalties, margin drag, and assignment-risk flags can erase the apparent Sharpe. This is a
useful publication result. It shows that covariance optimization and conditional option
premia are necessary for an option-only Markowitz paper, but they are not sufficient for a
tradable strategy claim.

\begin{figure}[t]
\centering
\includegraphics[width=0.88\textwidth]{portfolio_growth.pdf}
\caption{Growth of one dollar for option strategies and equity benchmarks. The log-scale
plot floors paths at $10^{-4}$ so every generated series remains visible.}
\label{fig:growth}
\end{figure}

\begin{figure}[t]
\centering
\includegraphics[width=0.70\textwidth]{random_sharpe_histogram.pdf}
\caption{Distribution of out-of-sample Sharpe ratios for random gross-NAV feasible option
portfolios, with the optimized option book marked.}
\label{fig:random}
\end{figure}

\subsection{Risk, attribution, and equity-drift controls}

\begin{table}[H]
\centering
\resizebox{0.92\textwidth}{!}{\begin{tabular}{lrrrr}
\toprule
Strategy & Predicted monthly vol & Train realized vol & Test realized vol & Test / predicted \\
\midrule
Equity-option Greek Markowitz & 0.334 & 0.267 & 0.307 & 0.918 \\
Greek Markowitz + VIX & 0.363 & 0.261 & 0.262 & 0.722 \\
Beta/delta-neutral + VIX & 0.354 & 0.272 & 0.239 & 0.677 \\
Equal premium & 1.509 & 0.625 & 0.542 & 0.359 \\
Equal risk & 0.668 & 0.394 & 0.368 & 0.551 \\
VIX hedge sleeve & 4.648 & 4.092 & 0.417 & 0.090 \\
\bottomrule
\end{tabular}
}
\caption{Greek-induced covariance calibration. Predicted monthly volatility is computed
from $q^\top\Sigma_Oq$ using the training-window Greek risk model; realized volatility is
measured on train and test bucket returns.}
\label{tab:riskcal}
\end{table}

\begin{figure}[t]
\centering
\includegraphics[width=0.66\textwidth]{risk_calibration.pdf}
\caption{Predicted versus realized monthly volatility by option strategy.}
\label{fig:riskcal}
\end{figure}

\begin{table}[H]
\centering
\resizebox{\textwidth}{!}{\begin{tabular}{lrrrrrrrrrrrrrrrrrrr}
\toprule
Strategy & Ann. alpha & Alpha HAC t & Alpha HAC se & HAC lags & $R^2$ & Residual ann. vol & N & Beta SPY & Beta AAPL & Beta AMZN & Beta GOOGL & Beta JPM & Beta META & Beta MSFT & Beta NVDA & Beta TSLA & Beta VX front & Beta dVIX & Beta dVVIX \\
\midrule
Equity-option Greek Markowitz & 0.405 & 1.141 & 0.355 & 3 & 0.618 & 0.634 & 58 & 0.005 & -1.331 & -0.875 & 0.631 & 0.590 & 0.126 & 2.585 & -0.089 & 1.136 & -1.368 & 0.046 & 0.003 \\
Greek Markowitz + VIX & 0.675 & 2.060 & 0.328 & 3 & 0.513 & 0.616 & 58 & 1.908 & -1.462 & -0.747 & 0.625 & -0.127 & -0.026 & 1.877 & 0.222 & 0.871 & -0.767 & 0.056 & -0.001 \\
Beta/delta-neutral + VIX & 0.845 & 2.753 & 0.307 & 3 & 0.429 & 0.601 & 58 & 0.715 & -1.475 & -0.612 & 0.472 & -0.579 & -0.170 & 1.336 & 0.370 & 0.815 & -0.522 & 0.038 & -0.001 \\
Equal premium & -0.416 & -0.566 & 0.735 & 3 & 0.253 & 1.647 & 58 & 2.522 & -1.714 & -0.069 & 1.943 & -1.879 & -0.344 & 3.001 & -0.374 & 0.404 & -0.749 & 0.032 & -0.001 \\
Equal risk & -0.183 & -0.334 & 0.547 & 3 & 0.253 & 1.115 & 58 & 2.041 & -1.795 & -0.113 & 1.331 & -1.484 & -0.316 & 1.968 & -0.082 & 0.321 & -0.379 & 0.020 & -0.003 \\
VIX hedge sleeve & -8.807 & -14.623 & 0.602 & 3 & 0.252 & 1.264 & 58 & 3.562 & -0.597 & -0.589 & 0.670 & -3.047 & 0.423 & -0.958 & -0.030 & -0.131 & -1.180 & 0.093 & -0.013 \\
Delta-matched equities & 0.000 & 2.229 & 0.000 & 3 & 1.000 & 0.000 & 58 & -0.000 & 0.088 & 0.539 & 0.230 & 0.590 & 0.011 & 0.442 & 0.779 & 0.628 & -0.000 & -0.000 & 0.000 \\
Underlying Markowitz & 0.000 & 2.700 & 0.000 & 3 & 1.000 & 0.000 & 58 & -0.000 & -0.085 & 0.159 & -0.096 & -0.050 & 0.029 & 0.280 & 0.191 & 0.111 & -0.000 & -0.000 & 0.000 \\
\bottomrule
\end{tabular}
}
\caption{Out-of-sample factor regressions on SPY, the eight underlying equities, VX front
returns, VIX changes, and VVIX changes. Alpha is annualized and its t-statistic uses
HAC/Newey--West standard errors.}
\label{tab:regression}
\end{table}

The regression is an audit, not a causal proof. If the option book is only a levered
large-cap equity book, SPY and single-name betas should explain it. If it is only a
volatility bet, VX, VIX, and VVIX controls should reveal that. The paper therefore reports
alpha, beta, residual volatility, and $R^2$ for every strategy and benchmark. A high point
estimate without a credible HAC t-statistic or without robustness to costs is not treated
as an economic discovery.

\begin{table}[H]
\centering
\resizebox{\textwidth}{!}{\begin{tabular}{lrrrrrrrrrrrr}
\toprule
Strategy & Equity delta & Equity gamma & Equity vega & VIX-forward delta & VIX-forward gamma & VIX-option vega & Theta/carry & VX roll & Skew/tail & Residual & Realized ann. mean & Residual vol share \\
\midrule
Equity-option Greek Markowitz & 0.521 & 0.898 & 0.002 & 0.000 & 0.000 & 0.000 & -0.003 & 0.000 & 0.000 & -0.524 & 0.894 & 0.570 \\
Greek Markowitz + VIX & 0.489 & 1.475 & 0.001 & 0.767 & -0.867 & -0.227 & 1.057 & 0.000 & 0.000 & -1.447 & 1.246 & 0.950 \\
Beta/delta-neutral + VIX & 0.168 & 1.486 & 0.001 & 0.554 & -0.745 & -0.190 & 1.034 & 0.000 & 0.000 & -1.135 & 1.172 & 1.069 \\
Cost-aware Sortino + VIX & 2.064 & 8.577 & 0.004 & 0.285 & -1.134 & -0.112 & 2.169 & 0.000 & 0.000 & -6.833 & 5.020 & 0.710 \\
Equal premium & 0.062 & 7.928 & 0.006 & -0.016 & 0.358 & 0.045 & -1.016 & 0.000 & 0.000 & -7.423 & -0.056 & 0.354 \\
Equal risk & 0.027 & 5.625 & 0.004 & -0.096 & 0.270 & 0.049 & -0.592 & 0.000 & 0.000 & -5.184 & 0.104 & 0.514 \\
VIX hedge sleeve & 0.000 & 0.000 & 0.000 & 2.083 & 2.615 & 0.135 & -12.254 & 0.000 & 0.000 & -1.865 & -9.285 & 4.858 \\
\bottomrule
\end{tabular}
}
\caption{Annualized mean P\&L attribution for the option-only strategies. Equity-option
Greeks and VIX-forward Greeks are separated; components sum to realized mean return after
the residual term.}
\label{tab:attribution}
\end{table}

\subsection{Regimes and robustness}

\begin{table}[H]
\centering
\resizebox{\textwidth}{!}{\begin{tabular}{llrrrrrrrrrr}
\toprule
Strategy & Regime & Months & Ann. return & Ann. vol & Sharpe & Sharpe CI lo & Sharpe CI hi & Sortino & Calmar & Omega & Info. ratio \\
\midrule
Equity-option Greek Markowitz & Down & 19 & -0.683 & 0.847 & -0.806 & -1.975 & 0.104 & -1.017 & -0.811 & 0.572 & -0.225 \\
Equity-option Greek Markowitz & Flat & 21 & 0.518 & 0.928 & 0.558 & -0.761 & 1.556 & 1.137 & 0.823 & 1.525 & -0.335 \\
Equity-option Greek Markowitz & Up & 20 & 2.787 & 1.175 & 2.373 & 0.685 & 4.553 & 8.069 & 6.615 & 6.928 & 1.808 \\
Greek Markowitz + VIX & Down & 19 & 0.806 & 0.746 & 1.081 & 0.153 & 1.938 & 2.583 & 1.918 & 2.235 & 1.764 \\
Greek Markowitz + VIX & Flat & 21 & 0.562 & 0.917 & 0.612 & -0.991 & 1.786 & 1.343 & 0.965 & 1.618 & 0.073 \\
Greek Markowitz + VIX & Up & 20 & 2.384 & 0.985 & 2.421 & 0.831 & 4.232 & 6.598 & 7.182 & 5.575 & 1.737 \\
Beta/delta-neutral + VIX & Down & 19 & 1.487 & 0.721 & 2.061 & 1.366 & 3.248 & 4.939 & 5.123 & 4.422 & 2.700 \\
Beta/delta-neutral + VIX & Flat & 21 & 0.649 & 0.815 & 0.796 & -0.666 & 1.860 & 2.154 & 1.757 & 1.990 & 0.151 \\
Beta/delta-neutral + VIX & Up & 20 & 1.423 & 0.951 & 1.496 & -0.182 & 3.218 & 3.419 & 3.357 & 3.019 & 0.777 \\
Equal premium & Down & 19 & -0.515 & 1.890 & -0.273 & -1.562 & 0.450 & -0.473 & -0.535 & 0.812 & -0.013 \\
Equal premium & Flat & 21 & -2.368 & 0.879 & -2.694 & -5.258 & -1.256 & -2.270 & -2.377 & 0.155 & -3.484 \\
Equal premium & Up & 20 & 2.809 & 2.352 & 1.195 & 0.314 & 2.010 & 3.236 & 3.712 & 2.597 & 0.917 \\
Equal risk & Down & 19 & 0.101 & 1.436 & 0.070 & -0.954 & 0.868 & 0.128 & 0.130 & 1.054 & 0.411 \\
Equal risk & Flat & 21 & -1.654 & 0.687 & -2.409 & -4.471 & -0.990 & -2.113 & -1.703 & 0.155 & -3.447 \\
Equal risk & Up & 20 & 1.954 & 1.419 & 1.377 & 0.302 & 2.376 & 3.665 & 3.495 & 2.843 & 0.915 \\
VIX hedge sleeve & Down & 19 & -7.743 & 1.368 & -5.662 & -9.800 & -4.100 & -2.982 &  & 0.021 & -5.203 \\
VIX hedge sleeve & Flat & 21 & -8.367 & 1.909 & -4.383 & -6.184 & -3.226 & -2.829 &  & 0.104 & -4.129 \\
VIX hedge sleeve & Up & 20 & -11.714 & 0.276 & -42.492 & -107.213 & -26.949 & -3.453 &  & 0.000 & -40.570 \\
Delta-matched equities & Down & 19 & -1.767 & 0.675 & -2.619 & -4.396 & -1.375 & -2.204 & -1.817 & 0.159 & -2.086 \\
Delta-matched equities & Flat & 21 & 1.272 & 0.707 & 1.800 & 0.747 & 2.788 & 5.921 & 3.839 & 5.091 & 1.172 \\
Delta-matched equities & Up & 20 & 3.630 & 0.698 & 5.203 & 3.511 & 8.613 & 34.716 & 26.891 & 45.819 & 4.428 \\
Underlying Markowitz & Down & 19 & -0.319 & 0.120 & -2.664 & -4.820 & -1.265 & -2.262 & -0.765 & 0.148 & 1.948 \\
Underlying Markowitz & Flat & 21 & 0.172 & 0.185 & 0.933 & -0.496 & 2.168 & 2.087 & 1.120 & 2.038 & -0.742 \\
Underlying Markowitz & Up & 20 & 0.658 & 0.172 & 3.829 & 2.463 & 6.121 & 15.942 & 17.455 & 15.149 & -0.128 \\
\bottomrule
\end{tabular}
}
\caption{Performance by SPY-return tercile. Down, flat, and up regimes are defined using
the test-window distribution of monthly SPY returns; Sharpe intervals use circular block
bootstrap resampling.}
\label{tab:regime}
\end{table}

\begin{figure}[t]
\centering
\includegraphics[width=0.88\textwidth]{regime_sharpes.pdf}
\caption{Sharpe ratios by SPY market regime. Markers and line styles are shown for every
reported strategy and benchmark.}
\label{fig:regime}
\end{figure}

The robustness message is intentionally conservative. VIX options are not treated as a
free return source; they are volatility derivatives whose settlement must be exact before
headline claims are allowed. Cost, capacity, margin, and assignment diagnostics determine
whether the gross research edge survives implementation. Appendix tables report the full
trading-data audit, exposure diagnostics, VIX-regime splits, leave-one-out tests, rolling
OOS check, and full inference ledger.
